Durante el período 2018–2025 se observa que el rendimiento acumulado de las empresas presenta una distribución heterogénea. NVIDIA registró el mayor crecimiento, alcanzando un rendimiento acumulado de
1131.46%
, seguida por Tesla con
990.14%
y BioNTech (BNTX) con
617.53%
. Estos resultados reflejan una apreciación significativa en el precio de sus acciones durante el período analizado, posicionándolas como las empresas con mayor retorno para un inversionista que hubiera mantenido su inversión desde el inicio hasta el final del horizonte de estudio.
Las empresas ubicadas entre la cuarta y la décima posición también presentan rendimientos acumulados positivos; sin embargo, muestran un comportamiento más homogéneo, con valores comprendidos entre
527.31%
y
478.65%
. Esta diferencia respecto a las tres primeras posiciones evidencia que el crecimiento del mercado no fue uniforme y que únicamente un grupo reducido de compañías logró obtener retornos excepcionalmente superiores al promedio.
Desde una perspectiva financiera, los resultados sugieren que las empresas del sector tecnológico dominaron el rendimiento bursátil durante el período analizado, consolidándose como uno de los segmentos con mayor generación de valor para los inversionistas. No obstante, el rendimiento histórico constituye únicamente una referencia y no garantiza resultados futuros, por lo que cualquier decisión de inversión debería complementarse con análisis fundamentales, evaluación del riesgo, condiciones del mercado y estudios estadísticos que permitan valorar la sostenibilidad del crecimiento observado.